% Complete Section 3.1, original PDF 31-32 / printed 21-22.
% block: unbounded-definition-heading page:31 type:section
\subsection{{无界算子}}\label{sec:unbounded-operators}

% block: unbounded-domain-definition pages:31,32 type:ordinary-paragraph
{考虑无界算子时, 必须允许算子的定义域不是整个 Hilbert 空间, 而只是一个子空间. 因此, 所谓从 Hilbert 空间 $X$ 到 Hilbert 空间 $Y$ 的\emph{无界算子} $T$, 是指将 $X$ 的一个子空间 $D(T)$ 映入 $Y$ 的线性算子, 其中该子空间称为 $T$ 的定义域. 算子 $T$ 不一定有界, 定义域 $D(T)$ 也不一定是闭子空间. 这类算子称为 Hilbert 空间上的无界算子, 尽管其中也包含在整个 $X$ 上定义的有界线性算子. 更准确的名称\emph{不一定处处有定义且不一定有界的线性算子}就有些冗长了! 我们几乎总是关注\emph{稠密定义}的算子, 即定义域 $D(T)$ 是 $X$ 的稠密子空间. 一个值得牢记的典型例子是: 对某个区域 $\Omega\subset\mathbb R^n$, 取 $X=L^2(\Omega)$, $Y=L^2(\Omega;\mathbb R^n)$, 定义域 $D(T)$ 为 Sobolev 空间 $H^1(\Omega)$, 它是 $L^2(\Omega)$ 的稠密子空间, 而算子 $T$ 为梯度. 也就是说, 对所有 $\phi\in H^1(\Omega)$, 有 $T\phi=\operatorname{grad}\phi\in L^2(\Omega;\mathbb R^n)$. 一个变体是算子 $\widetilde T$, 其定义域 $D(\widetilde T)=\mathring H^1(\Omega)$ 同样在 $L^2$ 中稠密, 且仍有 $\widetilde T\phi=\operatorname{grad}\phi$, 但 $\widetilde T\phi$ 只对 $\phi\in\mathring H^1$ 定义. 还可以定义其他梯度算子, 其定义域由 $\overline\Omega$ 上的 $C^1$ 函数, $C^\infty$ 函数, 或在 $\Omega$ 内有紧支集的 $C^\infty$ 函数组成. 尽管我们可能将这些算子中的任何一个都称为梯度, 它们却并不相同, 因而始终明确其定义域十分重要. 若说两个无界算子 $S$ 与 $T$ 相等, 则既是指 $D(S)=D(T)$, 也是指两算子的作用相同.}

% block: unbounded-sum-composition page:32 type:ordinary-paragraph
{从 $X$ 到 $Y$ 的两个无界算子 $S$ 和 $T$ 可以相加, 得到定义域为 $D(S+T)=D(S)\cap D(T)$ 的算子 $S+T$. 若 $S$ 是无界算子 $X\to Y$, $T$ 是无界算子 $Y\to Z$, 则类似地可以定义复合 $T\circ S$, 其定义域为 $\{\,v\in D(S)\mid Sv\in D(T)\,\}$. 必须注意, 稠密定义算子的和或复合不一定稠密定义.}

% block: null-range-graph-intro page:32 type:ordinary-paragraph
{无界算子 $T$ 的\emph{零空间}, $T$ 的\emph{像空间}和 $T$ 的\emph{图}分别按显然的方式定义:}

% block: null-range-graph-definitions page:32 type:equation
\[
\begin{aligned}
\mathcal N(T)&=\{\,x\in D(T)\mid Tx=0\,\}\subset X,\\
\mathcal R(T)&=\{\,Tx\mid x\in D(T)\,\}\subset Y,\\
\Gamma(T)&=\{\,(x,Tx)\mid x\in D(T)\,\}\subset X\times Y.
\end{aligned}
\]

% block: graph-norm-intro page:32 type:ordinary-paragraph
{函数 $v\mapsto(v,Tv)$ 是从 $T$ 的定义域到其图的一一映射. 将 $X\times Y$ 中的范数应用于 $(v,Tv)$, 就在 $D(T)$ 上定义了一个范数. 这个范数称为 $D(T)$ 上的\emph{图范数}, 由下式给出:}

% block: graph-norm-equation page:32 type:equation
\[
\|v\|_{D(T)}^2:=\|v\|_X^2+\|Tv\|_Y^2,\qquad v\in D(T),
\]

% block: graph-inner-product-intro page:32 type:ordinary-paragraph
{它与下述\emph{图内积}相联系:}

% block: graph-inner-product-equation page:32 type:equation
\[
\langle v,w\rangle_{D(T)}:=\langle v,w\rangle_X+\langle Tv,Tw\rangle_Y,\qquad v,w\in D(T).
\]

% block: graph-norm-strength page:32 type:ordinary-paragraph
{显然, 图范数强于 $X$ 范数在 $T(D)$ 上的限制, 而且通常两者不等价. 例如, 在前面的梯度算子例子中, 图范数就是通常的 $H^1$ 范数, 它严格强于 $L^2$ 范数. 总有}

% block: graph-boundedness-equation page:32 type:equation
\[
\|Tv\|_Y\leq\|v\|_{D(T)},\qquad v\in D(T),
\]

% block: graph-norm-completeness page:32 type:ordinary-paragraph
{因此, 当前者配备图范数时, 算子 $T$ 是从 $D(T)$ 到 $Y$ 的有界算子. 图范数使 $D(T)$ 成为内积空间, 或预 Hilbert 空间, 但它未必完备.}%
% block: graph-norm-footnote page:32 type:footnote
\footnote{{也可以不取平方, 将图范数定义为 $\|v\|_{D(T)}=\|v\|_X+\|Tv\|_Y$. 这个定义给出一个等价范数, 但它将不再与内积相联系.}}
